(1) Laplace方程轴对称且外驱动只含一阶角谐波，故内势可写 $\Phi_{\rm in}=-A r\cos\theta$，外势 $\Phi_{\rm out}=-E_0r\cos\theta+B\cos\theta/r^2$。不允许内势有 $r^{-2}$ 奇性，也不允许外势有与外场不符的增长项。

(2) 在 $r=a$ 处 $-Aa=-E_0a+B/a^2$，法向位移连续给 $\varepsilon_1A=\varepsilon_2(E_0+2B/a^3)$。联立得 $\boxed{A=3\varepsilon_2E_0/(\varepsilon_1+2\varepsilon_2)}$、$B=a^3E_0(\varepsilon_1-\varepsilon_2)/(\varepsilon_1+2\varepsilon_2)$。偶极势系数为 $p/(4\pi\varepsilon_2)=B$，故 $\mathbf p=4\pi\varepsilon_2a^3(\varepsilon_1-\varepsilon_2)E_0/(\varepsilon_1+2\varepsilon_2)\mathbf e_z$。

(3) 若外部介质不极化且内球极化率为 $\chi=\varepsilon_1/\varepsilon_2-1$，界面相对背景的极化面电荷可用 $\sigma_b=(\mathbf P_1-\mathbf P_2)\cdot\hat r$ 表示；对球内均匀极化和外背景已被外场极化的情况，等效源的面电荷是 $\boxed{\sigma_{\rm eff}=3\varepsilon_2(\varepsilon_1-\varepsilon_2)E_0\cos\theta/(\varepsilon_1+2\varepsilon_2)}$，积分 $\int_0^{2\pi}d\phi\int_0^\pi\cos\theta\sin\theta d\theta=0$，所以净束缚源电荷为零。

(4) 当 $\varepsilon_1=\varepsilon_2$，有 $A=E_0,B=0$，无扰动；当 $\varepsilon_1\gg\varepsilon_2$，$A\to0$，球近似导体，偶极矩趋 $4\pi\varepsilon_2a^3E_0$；当 $\varepsilon_1\to0^+$，$A\to3E_0/2$，内电场增强、诱导偶极与外场反向。
